\section{The fixed-\texorpdfstring{\(\omega\)}{omega} extremal problem}
\label{s-fixed-slice}

We normalize Robin's inequality by
\[
    G(n)
    =
    \frac{\sigma(n)}
    {e^\gamma n\log\log n}.
\]
Thus Robin's inequality is \(G(n)<1\) for \(n\geq5041\)
\cite{robin1984}. For \(m\geq1\), define the fixed-\(\omega\) extremal
value
\[
    g_{\integers}(m)
    =
    \sup_{\substack{n\geq5041\\ \omega(n)=m}}\log G(n).
\]
A negative value of \(g_{\integers}(m)\) means that Robin's inequality
holds throughout the slice. For \(m\geq6\), the attainment proved below
makes the converse true as well.

We write \(p_i\) for the \(i\)th prime, and use
\[
    \pi(x)=\#\{p\leq x\},
    \qquad
    \theta(x)=\sum_{p\leq x}\log p.
\]
We use \(\integers_{++}=\{1,2,\ldots\}\).
The first reduction removes the choice of prime support.

\begin{proposition}[Support reduction]
\label{prop-support}
For every \(m\geq6\), the supremum defining
\(g_{\integers}(m)\) is a maximum. Every maximizer has the form
\[
    n=\prod_{i=1}^{m}p_i^{r_i},
    \qquad
    r\in\integers_{++}^{m}.
\]
\end{proposition}

\begin{proof}
Suppose
\[
    n=\prod_{i=1}^{m}q_i^{r_i},
    \qquad
    q_1<\cdots<q_m,
    \qquad
    r_i\geq1.
\]
For a fixed exponent \(a\), the local factor
\[
    1+q^{-1}+\cdots+q^{-a}
\]
is strictly decreasing in \(q\). We replace \(q_i\) by \(p_i\) for every
\(i\). This cannot decrease \(\sigma(n)/n\) or increase \(n\). Since
\(\log\log n\) is increasing on the range in question, the replacement
cannot decrease \(G(n)\).

For \(m\geq6\), the replacement remains feasible because
\[
    \prod_{i=1}^{m}p_i
    \geq
    \prod_{i=1}^{6}p_i
    =
    30030
    >
    5041.
\]
If some \(q_i>p_i\), the corresponding local factor increases strictly,
while the denominator does not increase. Hence \(G(n)\) increases
strictly, and every maximizer uses \(p_1,\ldots,p_m\).

We now prove attainment. On this fixed support,
\(\sigma(n)/n\) is bounded above by
\[
    \prod_{i=1}^{m}\frac{p_i}{p_i-1},
\]
whereas \(\log\log n\to\infty\) when any exponent tends to infinity.
Hence \(\log G(n)\to-\infty\) whenever an exponent tends to infinity.
Only finitely many exponent vectors can
lie above the value at \(r_i=1\), so the supremum is attained.
\end{proof}

For \(m\leq5\) the product of the first \(m\) primes is smaller than
\(5041\), so the replacement in the proof can leave the range
\(n\geq5041\), and the support need not reduce to the first \(m\)
primes. A finite verification settles these slices.

\begin{proposition}[Small slices]
\label{prop-small-slices}
For \(1\leq m\leq5\), the supremum defining \(g_{\integers}(m)\) is a
maximum, attained uniquely at
\[
    2^{13},\qquad
    2^{6}3^{4},\qquad
    2^{5}3^{2}5^{2},\qquad
    2^{5}3^{2}5\cdot7,\qquad
    2^{4}3^{2}5\cdot7\cdot11,
\]
respectively. Each maximum is negative, and
\(g_{\integers}(4)=\log G(10080)=-0.0142\ldots\) exceeds the other
four.
\end{proposition}

\begin{proof}
The local factor of \(\sigma(n)/n\) at a prime \(p\) is smaller than
\(p/(p-1)\), and \(p/(p-1)\) decreases in \(p\), so \(\omega(n)=m\) gives
\(\sigma(n)/n<\prod_{i=1}^{m}p_i/(p_i-1)\). Hence
\[
    \log G(n)
    <
    \sum_{i=1}^{m}\log\frac{p_i}{p_i-1}
    -\gamma-\log\log\log n.
\]
The bound decreases in \(n\), and at \(n=10^7\) it is already below the
maximum claimed for the corresponding \(m\). A verification over the
integers in \([5041,10^7]\) with \(\omega(n)=m\) completes the proof.
\end{proof}

Robin's inequality therefore holds on every slice with \(m\leq5\). The
values are not monotone in \(m\), since
\(g_{\integers}(5)=\log G(55440)=-0.0168\ldots\) falls below the
\(m=4\) value. Section~\ref{s-bounded-support} shows that \(10080\)
holds the bounded-support maximum for \(6\leq m\leq29\).

\paragraph{Real-exponent relaxation.}
For the rest of the paper, \(m\geq6\) and the support is
\((p_1,\ldots,p_m)\). For \(r\in[1,\infty)^m\), we write
\[
    S(r)=\sum_{i=1}^{m}r_i\log p_i.
\]
For \(r\in\integers_{++}^{m}\), we have
\[
    \frac{\sigma(\prod_i p_i^{r_i})}
    {\prod_i p_i^{r_i}}
    =
    \prod_{i=1}^{m}
    \frac{p_i-p_i^{-r_i}}{p_i-1},
\]
so the real-exponent relaxation is obtained by using the expression on
the right for every \(r\in[1,\infty)^m\). Its value is
\begin{equation}
\label{e-objective}
    g_{\reals}(m)
    =
    \sup_{r\in[1,\infty)^m}
    \left(
        \sum_{i=1}^{m}\log(p_i-p_i^{-r_i})
        -\gamma
        -\sum_{i=1}^{m}\log(p_i-1)
        -\log\log S(r)
    \right).
\end{equation}
This objective agrees with \(\log G(\prod_i p_i^{r_i})\) at every
positive integer exponent vector.
Since the integer exponent vectors form a subset of its feasible set,
\[
    g_{\reals}(m)\geq g_{\integers}(m).
\]

The integrality gap is
\[
    \Igap(m)=g_{\reals}(m)-g_{\integers}(m).
\]
We refer to \(g_{\reals}(m)\) as the relaxed excess and to
\(-g_{\integers}(m)\) as the fixed-slice Robin margin. A positive relaxed
excess means that the real-exponent relaxation violates Robin's bound. A
positive Robin margin means that every integer in the slice satisfies the
bound.

Thus, for every \(m\geq6\),
\[
    -g_{\integers}(m)
    =
    \Igap(m)-g_{\reals}(m).
\]
We estimate the two terms separately. The gap is unconditional and comes
from the exponent lattice. The
relaxed excess is a normalized Mertens error and carries the dependence on
the zeros of the zeta function.
